\documentclass[aspectratio=169]{beamer}
\usetheme{tmg}
\usepackage[bigfiles]{pdfbase}
\ExplSyntaxOn
\NewDocumentCommand\embedvideo{smm}{
  \group_begin:
  \leavevmode
  \tl_if_exist:cTF{file_\file_mdfive_hash:n{#3}}{
    \tl_set_eq:Nc\video{file_\file_mdfive_hash:n{#3}}
  }{
    \IfFileExists{#3}{}{\GenericError{}{File~`#3'~not~found}{}{}}
    \pbs_pdfobj:nnn{}{fstream}{{}{#3}}
    \pbs_pdfobj:nnn{}{dict}{
      /Type/Filespec/F~(#3)/UF~(#3)
      /EF~<</F~\pbs_pdflastobj:>>
    }
    \tl_set:Nx\video{\pbs_pdflastobj:}
    \tl_gset_eq:cN{file_\file_mdfive_hash:n{#3}}\video
  }
  %
  \pbs_pdfobj:nnn{}{dict}{
    /Type/RichMediaInstance/Subtype/Video
    /Asset~\video
    /Params~<</FlashVars (
      source=#3&
      skin=SkinOverAllNoFullNoCaption.swf&
      skinAutoHide=true&
      skinBackgroundColor=0x5F5F5F&
      skinBackgroundAlpha=0
    )>>
  }
  %
  \pbs_pdfobj:nnn{}{dict}{
    /Type/RichMediaConfiguration/Subtype/Video
    /Instances~[\pbs_pdflastobj:]
  }
  %
  \pbs_pdfobj:nnn{}{dict}{
    /Type/RichMediaContent
    /Assets~<<
      /Names~[(#3)~\video]
    >>
    /Configurations~[\pbs_pdflastobj:]
  }
  \tl_set:Nx\rmcontent{\pbs_pdflastobj:}
  %
  \pbs_pdfobj:nnn{}{dict}{
    /Activation~<<
      /Condition/\IfBooleanTF{#1}{PV}{XA}
      /Presentation~<</Style/Embedded>>
    >>
    /Deactivation~<</Condition/PI>>
  }
  %
  \hbox_set:Nn\l_tmpa_box{#2}
  \tl_set:Nx\l_box_wd_tl{\dim_use:N\box_wd:N\l_tmpa_box}
  \tl_set:Nx\l_box_ht_tl{\dim_use:N\box_ht:N\l_tmpa_box}
  \tl_set:Nx\l_box_dp_tl{\dim_use:N\box_dp:N\l_tmpa_box}
  \pbs_pdfxform:nnnnn{1}{1}{}{}{\l_tmpa_box}
  %
  \pbs_pdfannot:nnnn{\l_box_wd_tl}{\l_box_ht_tl}{\l_box_dp_tl}{
    /Subtype/RichMedia
    /BS~<</W~0/S/S>>
    /Contents~(embedded~video~file:#3)
    /NM~(rma:#3)
    /AP~<</N~\pbs_pdflastxform:>>
    /RichMediaSettings~\pbs_pdflastobj:
    /RichMediaContent~\rmcontent
  }
  \phantom{#2}
  \group_end:
}
\ExplSyntaxOff

\title{Plantilla de Beamer de The Math Guys}
\subtitle{Versión 1.0}
\author{Nombre del autor}
\date{Fecha de la presentación}
% Los roles son los siguientes:
% - Administrador (Solo MathLike puede tener este rol)
% - Miembro
% - Moderador (Solo moderadores del server pueden tener este rol)
% - Ayudante de [asignatura] (Solo ayudantes de esa asignatura pueden tener este rol)
% - Creador de contenido (Solo creadores de contenido pueden tener este rol)
\institute{Nombre del rol de The Math Guys}
\titlegraphic{\includegraphics[height=0.25\paperheight]{TheMathGuysLogo.png}}
\discord{The math guys}
% Si no tiene YouTube o GitHub, simplemente borre el comando correspondiente
\youtube{Canal de YouTube del presentador}
\github{GitHub del presentador}

\begin{document}
\begin{frame}
  \titlepage
\end{frame}
\begin{frame}
  \frametitle{Hola, mundo}
  \begin{itemize}
    \item Esto es una lista.
    \item Pero sin orden.
  \end{itemize}
  \begin{enumerate}
    \item Esto es una lista.
    \item Pero con orden.
  \end{enumerate}
\end{frame}
\begin{frame}
  \frametitle{Otra diapositiva}
  $$\int_{-\infty}^{\infty}e^{-x^2}dx=\sqrt{\pi}$$
  Arriba está la fórmula de la integral de Gauss. \textit{Esto está en cursiva}, \textbf{esto en negrita}, \texttt{esto en fuente monoespaciada}, \alert{y yo soy quien llama la atención}.
  Viene una ecuación en la misma línea: $e^{i\pi}+1=0$.
\end{frame}
\begin{frame}[fragile]
  \frametitle{Un código insertado}
  \begin{minted}{python}
def factorial(n):
    if n==0:
        return 1
    else:
        return n*factorial(n-1)
  \end{minted}
\end{frame}
\begin{frame}
  \frametitle{Una imagen insertada}
  \begin{figure}
    \centering
    \includegraphics[width=0.3\textwidth]{TheMathGuysLogo.png}
    \caption{El logo de The Math Guys}
  \end{figure}
\end{frame}
\end{document}